\documentclass{article}

\usepackage{tabularx}
\usepackage{booktabs}
\usepackage[normalem]{ulem} % for strikethrough

\title{Problem Statement and Goals\\\progname}

\author{\authname}

\date{}

\input{../Comments}
%% Common Parts

\newcommand{\progname}{Industrial Plant Designer} % PUT YOUR PROGRAM NAME HERE
\newcommand{\authname}{Team 18, Chemware Engineering
\\ Kishor Pandya
\\ Dhrumil Pandya
\\ Madhu Sivapragasam
\\ Rasik Pokharel} % AUTHOR NAMES                  

\usepackage{hyperref}
    \hypersetup{colorlinks=true, linkcolor=blue, citecolor=blue, filecolor=blue,
                urlcolor=blue, unicode=false}
    \urlstyle{same}
                                


\begin{document}

\maketitle

\begin{table}[hp]
\caption{Revision History} \label{TblRevisionHistory}
\begin{tabularx}{\textwidth}{llX}
\toprule
\textbf{Date} & \textbf{Developer(s)} & \textbf{Change}\\
\midrule
September 18th & Kishor Pandya & Added Goals\\
September 21st & Madhu Sivapragasam & Added stretch goals, challenge level and extras and my part of the appendix\\
September 23rd & Dhrumil Pandya & Added reflection\\
September 23rd & Kishor Pandya & Updated goals and added answer to appendix\\
September 23rd & Rasik Pokharel & Updated Problem Statement and Development Plan section 1,2,3\\
\textcolor{blue}{April 3rd} & \textcolor{blue}{Madhu Sivapragasam} & \textcolor{blue}{Updates for revision 1}\\
\bottomrule
\end{tabularx}
\end{table}

\section{Problem Statement}







\subsection{What is the Problem that you're trying to solve?}

The current software used by Dr. Mahalec's research team for modeling industrial plant operations relies heavily on manual input through Excel tables for configuring plant structures. This approach introduces inefficiencies and increases the potential for user error during the setup process. Additionally, the lack of a graphical user interface (GUI) and an integrated database limits the system's ability to manage network topology effectively. 
\subsection{Why is this problem important?}

Addressing this problem is crucial for improving the efficiency and accuracy of industrial plant modeling. Manual data entry through Excel tables can be both time-consuming and prone to human error. Introducing errors in this process can lead to inaccurate models. The absence of a graphical user interface (GUI) and an integrated database complicates this process by limiting the ability to visualize, manage, and optimize complex plant networks. By enhancing the modeling and visualization process, Dr. Mahalec's team can be more effective, reduce the risk of errors, and streamline decision-making. Ultimately, solving this problem would contribute to more reliable and efficient plant operations modeling and planning.
 
\subsection{Inputs and Outputs}

\subsection*{Inputs}
\begin{itemize}
    \item \textbf{Plant Operation Data:} Pre-recorded data related to various operations within the plant.
    \item \textbf{User Parameters:} Inputs provided by users to configure the plant model.
    \item \textbf{Existing Python Models:} Existing python models are used in the backend to simulate and compute model outputs from the GUI. Since the python models are provided, we are considering them as inputs.
    \item \textbf{Simulation Parameters:} Variables and conditions set by the user to run different scenarios or simulations of plant operations. For example timeframe of simulations
    \item \textbf{Manual Adjustments:} Direct inputs from users to modify or adjust specific aspects of the plant model or its operating conditions.
\end{itemize}

\subsection*{Outputs}
\begin{itemize}
    \item \textbf{Visualized Plant Model:} Graphical representation of the plant's operations, showing how different components interact in real-time or under specified conditions.
    \item \textbf{Simulation Results:} Outputs from the Python models displayed through the GUI, such as production rates, efficiency metrics, and potential bottlenecks.
    \item \textbf{Database Storage:} Storage of input data, simulation results, and user configurations in the database for future reference and analysis.
    \item \textbf{Data Export:} Exportable data formats (e.g., CSV, Excel) for further analysis or record-keeping outside the software.

\end{itemize}


\subsection{Stakeholders}

\begin{enumerate}
    \item \textbf{Students (Project Team):}
    \begin{itemize}
        \item As the primary developers, the student team is responsible for designing, developing, and implementing the software. Their learning, skill development, and successful project completion are directly impacted by this project.
    \end{itemize}

    \item \textbf{Professor Vladimir Mahalec:}
    \begin{itemize}
        \item As the project supervisor and client, Professor Mahalec provides guidance, feedback, and expertise. He ensures the project meets academic standards and aligns with the course's learning objectives. 
    \end{itemize}

    \item \textbf{McMaster University (Engineering Department):}
    \begin{itemize}
        \item The university and its engineering department have an interest in the project's successful completion to showcase the quality of the engineering program, the capstone course as well as the quality of the students.
    \end{itemize}

    \item \textbf{End Users (Plant Operators/Engineers):}
    \begin{itemize}
        \item If the software is to be used in a real-world chemical plant, the operators and engineers who will use the software are key stakeholders. They will require an intuitive user interface to model, optimize, and plan the plant operations effectively.
    \end{itemize}

    \begin{itemize}
        \item Although the project doesn't involve direct collaboration with the chemical, the organization is a stakeholder. They may be interested in applying the software for improved plant operations or decision-making in the future.
    \end{itemize}

    \item \textbf{Software Engineering Capstone Course Coordinators (Dr. Smith):}
    \begin{itemize}
        \item They oversee the capstone projects, ensuring they meet educational goals and providing support to students and faculty advisors. They are interested in the successful delivery and demonstration of the project.
    \end{itemize}

    \item \textbf{Future Students:}
    \begin{itemize}
        \item Future cohorts of students may use the documentation as a learning resource or as a basis for further development and enhancement in subsequent projects.
    \end{itemize}

\end{enumerate}

\subsection{Environment}

\subsection*{Hardware Environment}
\begin{itemize}
    \item \textbf{Development Machines:} Laptops or desktop computers with sufficient processing power, memory (RAM), and storage to support software development and testing. Recommended specifications include at least an Intel i5 processor, 8GB of RAM, and 25 gigabits of available Solid State Drive storage.
    \item \textbf{Servers :} Since the project requires one or more centralized databases, a back-end server, and a front-end server, a physical machine will be required with the capability to host all software components.
    \item \textbf{Testing Environment:} Test environments may include virtual machines or Docker containers to simulate different operating systems and configurations for compatibility testing.
    \item \textbf{Peripheral Devices:} Monitors, keyboards, and other peripherals to facilitate the development and testing for the project.
\end{itemize}

\subsection*{Software Environment}
\begin{itemize}
    \item \textbf{Operating System:} Development is likely to be conducted on Windows, macOS, or Linux, depending on the team members' preferences and the requirements of the existing Python scripts.
    \item \textbf{Programming Languages:} 
    \begin{itemize}
        \item \textbf{Python:} Used for interfacing with existing plant models and performing simulations.
        \item \textbf{JavaScript/TypeScript:} For front-end development using web-based frameworks like react.
        \item \textbf{SQL/NoSQL:} Using SQL or NoSQL Database is going to be essential to address the storage and data organization requirements for this project.
    \end{itemize}
    \item \textbf{Development Frameworks:}
    \begin{itemize}
        \item \textbf{Front-end GUI:} React will be leveraged for building the GUI.
        \item \textbf{Back-end:} A web framework such as django or express will be required to implement the server-side components for this project, to interface with the DB and the back-end scripts.
    \end{itemize}
    \item \textbf{Database Management System (DBMS):} SQLite, MySQL, or PostgreSQL to store and manage data for the administration of the application as well as the data generated by the plant model and user interactions.
    \item \textbf{IDE and Code Editors:} Integrated Development Environments (IDEs) such as PyCharm or Visual Studio Code for efficient code development.
    \item \textbf{Version Control System:} Git and GitHub will be leveraged for version control and collaboration among team members.
    \item \textbf{Testing Tools:} Unit testing frameworks may be leveraged for front-end testing to ensure software reliability and performance.
    \item \textbf{Virtualization Tools:} Docker or virtual environments will be leveraged for isolating and managing different software configurations and dependencies.
\end{itemize}
\section{Goals}

\subsection{Build a comprehensive factory model}
The software must enable users to construct a visual model of a factory or plant. Each machine or process milestone will be represented as a node, with connections between nodes formed through edges, indicating the direction of workflows. This feature provides users with a high-level, intuitive understanding of the factory's layout and processes, facilitating planning and optimization.

\subsection{\textcolor{blue}{Ensure version management system for models}}
\textcolor{blue}{\sout{The software must save the user-generated factory model to a local database. This allows for easy retrieval, modification, or analysis at any later stage, ensuring that work is never lost and can be iteratively refined.}}
\textcolor{blue}{The software must be able to act as a version control system for factory models allowing for retireval, modification and analysis of these models at a later stage ensuring that work is never lost and can be iteratively refined.}

\subsection{Support simulation of factory operations}
The software must allow users to input relevant data into the nodes and run a simulation. By integrating with a script provided by Professor Mahalec's team, the system will compute output variables and present results to the user. This feature enables users to test and optimize factory operations before implementation, reducing risks and improving efficiency.

\subsection{Develop a functional and intuitive GUI}
The graphical user interface (GUI) should allow users to:
\begin{itemize}
    \item Design and visualize a factory model using nodes and edges.
    \item Set time horizons for tracking or simulating production over specific periods.
    \item Display output from the simulation script in a clear, user-friendly manner.
\end{itemize}
The GUI's functionality will be validated through iterative testing, ensuring ease of use and efficiency.

\subsection{Emphasize code readability and testing}
The software will be built with clean, well-documented code, ensuring that it is easily understandable and maintainable by other engineers. A comprehensive test suite will be developed to ensure that all functionalities are thoroughly tested and robust. Code quality will be tracked through documentation reviews, and test coverage will be monitored to ensure the system's reliability.


\section{Stretch Goals}

\subsection{Develop a cloud deployed solution}
Time permitting we should try to deploy our service to the cloud so that a user can access the product from the web rather than having to deploy locally. This would increase the ease of use of our application since users would not have to complete any sort of local deployment steps and instead just access a specific URL. It would also allow our service to be used by more people since it would be accessible to everyone from the web. 

\subsection{Create a customizable user interface}
Since we're giving our sponsor full ownership of the product after it's been completed, it would be beneficial for them to have the ability to customize the user interface we create via some sort of admin tooling in case they want to change the UI at a later date. They should be able to do this with minimal coding knowledge and so a feature that enables them to customize different components could be useful in the transfer of ownership. This would improve the convenience of our application since the owners can customize it to fit their needs. 

\subsection{\textcolor{blue}{\sout{Enable multi account authentication}}}
\textcolor{blue}{\sout{As an extension of our version management software, it would be useful to add integrations for having multiple different accounts with authentication for each account. This is also a requirement to enable cloud deployments to prevent bad actors from accessing the service. Adding this feature would improve the security of our application while also making it more flexible by enabling multiple different parties to access and use the service in their own unique ways.} No longer a goal since its not in the scope of thsi project. }  


\section{Challenge Level and Extras}

The challenge level for our project is general. It is not basic because our project is a complex full-stack project that requires the use of both backend services to allow for version management while also building complex front-end components to enable the user to create and modify plant schemas. However, there is not enough novelty to be considered an advanced project since it doesn't have the complexity that requires considerable research or innovation. However, the plant design and version management features will require complex and difficult development so we think it easily meets the requirements for a general challenge level project. 

\textcolor{blue}{\sout{For extras, we plan on adding user documentation. Since our sponsor is taking ownership of the project after we are complete, they must have access to user documentation to properly use the service. Therefore adding user documentation is one suitable extra component that helps contribute to the overall challenge level of our project. Another extra we will add is usability testing. This is to ensure our project is user-centered and fits all the needs/wants set by our users. It also helps ensure our systems are robust and useable while also contributing to the overall challenge level of our project.} Since we are now handing over the project to a new set of developers at the end of the term, we decided that creating a code walkthrough  and a user instructional video made more sense as extras. This is because our existing docs provide enough text based content for the user to follow and it would be easier for newer developers to learn how to use the system through a video. Similarly we thought the code walkthrough would be useful as it allows newer developers to learn the system faster.} 

\newpage{}

\section*{Appendix --- Reflection}

\input{../Reflection.tex}

\begin{enumerate}
    \item What went well while writing this deliverable? 
    \\\\ Madhu: I found that our group did well in coordinating the work and splitting up tasks effectively. I think this is a good sign for future work since it means we should do well at delegating work throughout the project and thats a good habit to build early. We also all completed our work on time effectively which helps improve the overall trust we have in each other. 
    \\\\ Dhrumil: For me, one thing that went really well was the support we received from Professor Mahalec and his team. They were always available to clarify any questions we had about the scope of the project, which helped us align our understanding early on. Having their insights made a big difference, especially when we were defining our objectives and figuring out how to approach the more technical aspects of the project. It helped us feel confident that we were on the right track and reduced a lot of the uncertainty that can sometimes come with a project of this scope.
    \\\\ Kishor: While working on this deliverable, our team maintained a strong sense of collaboration and clear communication. We were able to break down the project into smaller sections, which made it easier for each member to focus on specific tasks. Additionally, we aligned our goals early on, which ensured that everyone was on the same page regarding the overall objectives.
     \\\\ Rasik: We worked well as a team, communicating effectively throughout using various platforms, including discord, teams, and SMS. All meetings held with dr Mahalac were concise and effective.
    \item What pain points did you experience during this deliverable, and how
    did you resolve them?
    \\\\ Madhu: We found it difficult to create design plans and problem statements with limited information on the project and so we had to schedule quick meetings with the sponsor to get a better understanding of the problem domain to write this document. To avoid the lack of information problem in the future, we've scheduled regular meetings with our sponsor going forward so that we can always get information when we need it. 
    \\\\ Dhrumil: One of the main pain points we faced was working with the predefined requirements from the supervisor, particularly around the use of a relational database. While we understand the reasoning behind this choice, we believe there might be a better fit for our specific use case when it comes to the user database. Given the nature of the project, a different database structure could potentially be more efficient. We haven’t completely ruled out the relational database option yet, but we are currently weighing the pros and cons of alternative solutions. Our plan is to compile these findings and present them to Professor Mahalec for further discussion. This way, we can ensure that we’re using the most appropriate tool for the project while still considering the original requirements.
    \\\\ Kishor: One challenge we faced was aligning everyone's understanding of the project requirements, particularly in relation to the technical aspects, such as the integration with external scripts. I initially found it difficult to grasp the level of detail required for certain sections, especially when defining goals in a way that avoided implementation specifics. To resolve this, we held additional meetings to clarify any misunderstandings and ensured that all team members had access to relevant resources, such as examples from similar projects. Another pain point was ensuring that our goals were measurable and formulated as strong selling points. We overcame this by reviewing the rubric carefully and cross-referencing our goals to ensure they met the criteria, revising them when necessary to focus more on the “what” rather than the “how.”
    \\\\ Rasik: One of the challenges we had during this project was incorporating all requirements into the documentation. We were constantly learning new details about the project and constantly made documentation updates. Using github's merging and reviewing, while effective, can be a tedious process for documentation review.
    \item How did you and your team adjust the scope of your goals to ensure
    they are suitable for a Capstone project (not overly ambitious but also of
    appropriate complexity for a senior design project)?
    \\\\We attempted to use our own previous experience with software development with the planned technologies in order to better understand the scope of our project. In some cases we had to make assumptions but we tried to back that up with information from our sponsor as well as our own research into the complexity of different tools. However we are still flexible if the scope changes or we find that certain components are more or less complex than we considered which should allow us to ensure that our project has an overall appropriate complexity. 
\end{enumerate}  

\end{document}
